% Options for packages loaded elsewhere
\PassOptionsToPackage{unicode}{hyperref}
\PassOptionsToPackage{hyphens}{url}
\PassOptionsToPackage{dvipsnames,svgnames,x11names}{xcolor}
\documentclass[
  11pt,
]{article}
\usepackage{xcolor}
\usepackage[margin=1in]{geometry}
\usepackage{amsmath,amssymb}
\setcounter{secnumdepth}{-\maxdimen} % remove section numbering
\usepackage{iftex}
\ifPDFTeX
  \usepackage[T1]{fontenc}
  \usepackage[utf8]{inputenc}
  \usepackage{textcomp} % provide euro and other symbols
\else % if luatex or xetex
  \usepackage{unicode-math} % this also loads fontspec
  \defaultfontfeatures{Scale=MatchLowercase}
  \defaultfontfeatures[\rmfamily]{Ligatures=TeX,Scale=1}
\fi
\usepackage{lmodern}
\ifPDFTeX\else
  % xetex/luatex font selection
\fi
% Use upquote if available, for straight quotes in verbatim environments
\IfFileExists{upquote.sty}{\usepackage{upquote}}{}
\IfFileExists{microtype.sty}{% use microtype if available
  \usepackage[]{microtype}
  \UseMicrotypeSet[protrusion]{basicmath} % disable protrusion for tt fonts
}{}
\makeatletter
\@ifundefined{KOMAClassName}{% if non-KOMA class
  \IfFileExists{parskip.sty}{%
    \usepackage{parskip}
  }{% else
    \setlength{\parindent}{0pt}
    \setlength{\parskip}{6pt plus 2pt minus 1pt}}
}{% if KOMA class
  \KOMAoptions{parskip=half}}
\makeatother
\usepackage{longtable,booktabs,array}
\newcounter{none} % for unnumbered tables
\usepackage{calc} % for calculating minipage widths
% Correct order of tables after \paragraph or \subparagraph
\usepackage{etoolbox}
\makeatletter
\patchcmd\longtable{\par}{\if@noskipsec\mbox{}\fi\par}{}{}
\makeatother
% Allow footnotes in longtable head/foot
\IfFileExists{footnotehyper.sty}{\usepackage{footnotehyper}}{\usepackage{footnote}}
\makesavenoteenv{longtable}
\setlength{\emergencystretch}{3em} % prevent overfull lines
\providecommand{\tightlist}{%
  \setlength{\itemsep}{0pt}\setlength{\parskip}{0pt}}
\usepackage{bookmark}
\IfFileExists{xurl.sty}{\usepackage{xurl}}{} % add URL line breaks if available
\urlstyle{same}
\hypersetup{
  colorlinks=true,
  linkcolor={Maroon},
  filecolor={Maroon},
  citecolor={Blue},
  urlcolor={Blue},
  pdfcreator={LaTeX via pandoc}}

\author{}
\date{}

\begin{document}

\section{Mixed Wet-Half Simulation --- Brief Meeting
Report}\label{mixed-wet-half-simulation-brief-meeting-report}

\subsection{Purpose}\label{purpose}

This note summarises the current mixed wet-half spiral-inlet BOC
separator simulation. The aim is to clearly separate:

\begin{enumerate}
\def\labelenumi{\arabic{enumi}.}
\tightlist
\item
  what was calculated,
\item
  what was set up and tested, and
\item
  what was found from the current results.
\end{enumerate}

The main clarification is that this case used the reported Purnanto
inlet velocity of \textbf{26.81 m/s}, applied to the current geometry.
The inlet velocity was \textbf{not recalculated using the actual
inlet-half area} to force the total mass flow to exactly match the
Purnanto 1600 kJ/kg case.

\begin{center}\rule{0.5\linewidth}{0.5pt}\end{center}

\subsection{1. What Was Calculated}\label{what-was-calculated}

\subsubsection{1.1 Reference Purnanto inlet
condition}\label{reference-purnanto-inlet-condition}

The Purnanto 1600 kJ/kg reference inlet values were used as the
comparison target.

{\def\LTcaptype{none} % do not increment counter
\begin{longtable}[]{@{}lr@{}}
\toprule\noalign{}
Quantity & Value \\
\midrule\noalign{}
\endhead
\bottomrule\noalign{}
\endlastfoot
Liquid mass flow & 116.92 kg/s \\
Steam mass flow & 80.69 kg/s \\
Total mass flow & 197.61 kg/s \\
Steam quality by mass & 0.4083 \\
\end{longtable}
}

The steam quality is calculated as:

\[
x = \frac{\dot{m}_{\text{steam}}}
{\dot{m}_{\text{liquid}} + \dot{m}_{\text{steam}}}
\]

\[
x = \frac{80.69}{116.92 + 80.69}
= 0.4083
\]

These values were treated as reference values, not values exactly
imposed in the current simulation.

\begin{center}\rule{0.5\linewidth}{0.5pt}\end{center}

\subsubsection{1.2 Phase mass-flow
calculation}\label{phase-mass-flow-calculation}

The phase mass flow through the inlet is calculated from:

\[
\dot{m}_{\text{phase}}
=
\alpha_{\text{phase}}
\rho_{\text{phase}}
V A
\]

where:

\[
\alpha_{\text{phase}} = \text{phase volume fraction}
\]

\[
\rho_{\text{phase}} = \text{phase density}
\]

\[
V = \text{inlet velocity}
\]

\[
A = \text{inlet area}
\]

For the current simulation:

\[
V = 26.81\ \text{m/s}
\]

The velocity was \textbf{not} recalculated as an actual-area-corrected
velocity. Therefore, the inlet mass flow is close to the Purnanto
target, but it does not match it exactly.

\begin{center}\rule{0.5\linewidth}{0.5pt}\end{center}

\subsubsection{1.3 Volume-fraction calculation used for the mixed
wet-half
inlet}\label{volume-fraction-calculation-used-for-the-mixed-wet-half-inlet}

Fluent volume fraction is based on phase volumetric flow, not phase mass
fraction.

Using the Purnanto 1600 kJ/kg reference phase mass flows and densities:

\[
Q_{\text{liquid}}
=
\frac{\dot{m}_{\text{liquid}}}{\rho_{\text{liquid}}}
=
\frac{116.92}{881.77}
=
0.13260\ \text{m}^3/\text{s}
\]

\[
Q_{\text{steam}}
=
\frac{\dot{m}_{\text{steam}}}{\rho_{\text{steam}}}
=
\frac{80.69}{5.73}
=
14.08202\ \text{m}^3/\text{s}
\]

\[
Q_{\text{total}}
=
0.13260 + 14.08202
=
14.21462\ \text{m}^3/\text{s}
\]

The bulk liquid volume fraction for the full inlet flow is:

\[
\alpha_{\text{liquid,bulk}}
=
\frac{Q_{\text{liquid}}}{Q_{\text{total}}}
=
\frac{0.13260}{14.21462}
=
0.009328
\]

Because the current setup uses an equal two-half inlet and places all
liquid only in the wet half, the wet-half liquid volume fraction is
doubled:

\[
\alpha_{\text{liquid,wet-half}}
=
2\alpha_{\text{liquid,bulk}}
=
2(0.009328)
=
0.018656
\]

\[
\alpha_{\text{steam,wet-half}}
=
1 - 0.018656
=
0.981344
\]

Therefore, the inlet volume fractions used for this case were:

{\def\LTcaptype{none} % do not increment counter
\begin{longtable}[]{@{}lrr@{}}
\toprule\noalign{}
Boundary & Liquid volume fraction & Steam volume fraction \\
\midrule\noalign{}
\endhead
\bottomrule\noalign{}
\endlastfoot
Steam-only inlet half & 0.0 & 1.0 \\
Wet inlet half & 0.018656 & 0.981344 \\
\end{longtable}
}

This volume-fraction calculation comes from the velocity-inlet setup. It
should not be described as an actual-area-corrected setup. The
actual-area report uses the same volume fractions but recalculates the
velocity to 27.12 m/s to force the exact Purnanto mass flow; that was
\textbf{not} the setup used in this current case.

\begin{center}\rule{0.5\linewidth}{0.5pt}\end{center}

\subsubsection{1.4 Current simulation inlet flux
comparison}\label{current-simulation-inlet-flux-comparison}

From the current Fluent flux report:

{\def\LTcaptype{none} % do not increment counter
\begin{longtable}[]{@{}lr@{}}
\toprule\noalign{}
Phase & Current inlet flux \\
\midrule\noalign{}
\endhead
\bottomrule\noalign{}
\endlastfoot
Liquid inlet flow & 115.52 kg/s \\
Steam inlet flow & 80.71 kg/s \\
Total inlet flow & 196.23 kg/s \\
\end{longtable}
}

The current total inlet mass flow is:

\[
\dot{m}_{\text{current,total}}
=
115.5161 + 80.7123
=
196.2284\ \text{kg/s}
\]

Compared with the Purnanto target:

\[
\dot{m}_{\text{Purnanto,total}}
=
197.61\ \text{kg/s}
\]

The difference is:

\[
\Delta \dot{m}
=
196.2284 - 197.61
=
-1.3816\ \text{kg/s}
\]

The percentage difference is:

\[
\%\ \text{difference}
=
\frac{-1.3816}{197.61}
\times 100
=
-0.70\%
\]

Therefore, the current inlet flow is approximately \textbf{0.70\% lower}
than the Purnanto target.

\begin{center}\rule{0.5\linewidth}{0.5pt}\end{center}

\subsubsection{1.5 Current outlet-based
indicators}\label{current-outlet-based-indicators}

From the current outlet flux report:

{\def\LTcaptype{none} % do not increment counter
\begin{longtable}[]{@{}lr@{}}
\toprule\noalign{}
Quantity & Value \\
\midrule\noalign{}
\endhead
\bottomrule\noalign{}
\endlastfoot
Steam outlet flow & 81.45 kg/s \\
Liquid outlet flow & 2.50 kg/s \\
\end{longtable}
}

The liquid separation indicator based only on the outlet flux is:

\[
\eta_{\text{liquid}}
=
\frac{\dot{m}_{\text{liquid,outlet}}}
{\dot{m}_{\text{liquid,inlet}}}
\]

\[
\eta_{\text{liquid}}
=
\frac{2.4986}{115.5161}
=
0.02163
=
2.16\%
\]

The steam outlet dryness is:

\[
x_{\text{outlet}}
=
\frac{\dot{m}_{\text{steam,outlet}}}
{\dot{m}_{\text{steam,outlet}} + \dot{m}_{\text{liquid,outlet}}}
\]

\[
x_{\text{outlet}}
=
\frac{81.4524}{81.4524 + 2.4986}
=
0.9702
=
97.02\%
\]

These outlet-based values are preliminary because they still need to be
interpreted together with contours, vectors, pathlines, residuals, and
liquid hold-up inside the separator.

\begin{center}\rule{0.5\linewidth}{0.5pt}\end{center}

\subsection{2. What Was Set Up and
Tested}\label{what-was-set-up-and-tested}

\subsubsection{2.1 CFD setup used for this
case}\label{cfd-setup-used-for-this-case}

{\def\LTcaptype{none} % do not increment counter
\begin{longtable}[]{@{}ll@{}}
\toprule\noalign{}
Item & Current setup \\
\midrule\noalign{}
\endhead
\bottomrule\noalign{}
\endlastfoot
Geometry & Spiral-inlet BOC separator \\
Inlet type & Mixed wet-half split inlet \\
Boundary condition & Velocity inlet on both inlet halves \\
Inlet velocity & 26.81 m/s \\
Wet-half liquid volume fraction & 0.018656 \\
Wet-half steam volume fraction & 0.981344 \\
Steam-only inlet liquid volume fraction & 0.0 \\
Liquid density & 881.77 kg/m³ \\
Steam density & 5.73 kg/m³ \\
\end{longtable}
}

The key point is that the simulation used a \textbf{Purnanto
velocity-based inlet condition applied to the current geometry}.

It was not an exact \textbf{actual-area-corrected mass-flow match}.

\begin{center}\rule{0.5\linewidth}{0.5pt}\end{center}

\subsubsection{2.2 DPM particle injection
setup}\label{dpm-particle-injection-setup}

DPM particle injections were based on the Purnanto droplet-size
mass-flow distribution. The nine injection bins sum to:

\[
\dot{m}_{\text{DPM,total}}
=
116.92\ \text{kg/s}
\]

The bin sum is:

\[
\begin{aligned}
&5.846 + 13.708474 + 15.722479 + 15.224784 + 7.52414 \\
&+ 19.7663 + 17.859486 + 17.182465 + 4.085872 \\
&= 116.92\ \text{kg/s}
\end{aligned}
\]

The injection sizes and mass-flow weights used were:

{\def\LTcaptype{none} % do not increment counter
\begin{longtable}[]{@{}rrr@{}}
\toprule\noalign{}
Injection & Droplet diameter & Purnanto mass-flow weight \\
\midrule\noalign{}
\endhead
\bottomrule\noalign{}
\endlastfoot
1 & 1.29E-04 m & 5.846 kg/s \\
2 & 2.15E-04 m & 13.708474 kg/s \\
3 & 3.02E-04 m & 15.722479 kg/s \\
4 & 3.88E-04 m & 15.224784 kg/s \\
5 & 4.31E-04 m & 7.52414 kg/s \\
6 & 5.60E-04 m & 19.7663 kg/s \\
7 & 7.32E-04 m & 17.859486 kg/s \\
8 & 9.91E-04 m & 17.182465 kg/s \\
9 & 1.25E-03 m & 4.085872 kg/s \\
& \textbf{Total} & \textbf{116.92 kg/s} \\
\end{longtable}
}

These original Purnanto mass-flow weights were used rather than
actual-area-scaled values. This keeps the DPM comparison aligned with
the Purnanto 1600 kJ/kg liquid loading.

However, the DPM total is slightly higher than the current Eulerian
liquid inlet flux:

\[
\dot{m}_{\text{liquid,current}}
=
115.5161\ \text{kg/s}
\]

If the DPM bins need to be strictly scaled to the current inlet flux,
the scaling factor would be:

\[
\text{scale factor}
=
\frac{\dot{m}_{\text{liquid,current}}}
{\dot{m}_{\text{liquid,Purnanto}}}
\]

\[
\text{scale factor}
=
\frac{115.5161}{116.92}
=
0.98799
\]

Each Purnanto DPM bin could then be scaled as:

\[
\dot{m}_{\text{bin,scaled}}
=
\dot{m}_{\text{bin,Purnanto}}
\times 0.98799
\]

\begin{center}\rule{0.5\linewidth}{0.5pt}\end{center}

\subsubsection{2.3 DPM step sensitivity
check}\label{dpm-step-sensitivity-check}

The step sensitivity check used:

\begin{itemize}
\tightlist
\item
  injection tested: injection 1 / particle size 1,
\item
  initial setup: random initial setup,
\item
  particle streams: 100,
\item
  stochastic tracking: enabled,
\item
  eddy interaction/effect: enabled,
\item
  number of tries: 5,
\item
  effective total stochastic tracks: 500.
\end{itemize}

{\def\LTcaptype{none} % do not increment counter
\begin{longtable}[]{@{}
  >{\raggedleft\arraybackslash}p{(\linewidth - 10\tabcolsep) * \real{0.1667}}
  >{\raggedleft\arraybackslash}p{(\linewidth - 10\tabcolsep) * \real{0.1667}}
  >{\raggedleft\arraybackslash}p{(\linewidth - 10\tabcolsep) * \real{0.1667}}
  >{\raggedleft\arraybackslash}p{(\linewidth - 10\tabcolsep) * \real{0.1667}}
  >{\raggedleft\arraybackslash}p{(\linewidth - 10\tabcolsep) * \real{0.1667}}
  >{\raggedleft\arraybackslash}p{(\linewidth - 10\tabcolsep) * \real{0.1667}}@{}}
\toprule\noalign{}
\begin{minipage}[b]{\linewidth}\raggedleft
Max steps
\end{minipage} & \begin{minipage}[b]{\linewidth}\raggedleft
Step factor
\end{minipage} & \begin{minipage}[b]{\linewidth}\raggedleft
DPM iteration interval
\end{minipage} & \begin{minipage}[b]{\linewidth}\raggedleft
Trapped
\end{minipage} & \begin{minipage}[b]{\linewidth}\raggedleft
Incomplete
\end{minipage} & \begin{minipage}[b]{\linewidth}\raggedleft
Escaped
\end{minipage} \\
\midrule\noalign{}
\endhead
\bottomrule\noalign{}
\endlastfoot
50,000 & 5 & 10 & 158 & 342 & 0 \\
500,000 & 5 & 10 & 158 & 342 & 0 \\
50,000 & 2 & 2 & 159 & 341 & 0 \\
50,000 & 1 & 1 & 157 & 343 & 0 \\
500,000 & 2 & 2 & 159 & 341 & 0 \\
\end{longtable}
}

Increasing max steps from 50,000 to 500,000 did not reduce incomplete
tracks for this tested particle case. Reducing step factor and DPM
iteration interval produced only negligible changes in trapped and
incomplete counts. The incomplete-track issue therefore does not appear
to be solved simply by increasing max tracking steps for this setup.

The provisional future DPM setting from this check was:

{\def\LTcaptype{none} % do not increment counter
\begin{longtable}[]{@{}lr@{}}
\toprule\noalign{}
DPM setting & Value \\
\midrule\noalign{}
\endhead
\bottomrule\noalign{}
\endlastfoot
Max steps & 50,000 \\
Step factor & 2 \\
DPM iteration interval & 2 \\
Particle streams & 100 \\
Tries & 5 \\
Effective stochastic tracks & 500 per injection \\
\end{longtable}
}

\begin{center}\rule{0.5\linewidth}{0.5pt}\end{center}

\subsubsection{2.4 DPM stream-count sensitivity
check}\label{dpm-stream-count-sensitivity-check}

The stream-count check used the following fixed settings:

{\def\LTcaptype{none} % do not increment counter
\begin{longtable}[]{@{}lr@{}}
\toprule\noalign{}
Fixed setting & Value \\
\midrule\noalign{}
\endhead
\bottomrule\noalign{}
\endlastfoot
Max steps & 50,000 \\
Step factor & 2 \\
DPM iteration interval & 2 \\
Tries & 5 \\
Stochastic tracking tries & 5 \\
Random eddy lifetime & On \\
\end{longtable}
}

Stream-count results were:

{\def\LTcaptype{none} % do not increment counter
\begin{longtable}[]{@{}
  >{\raggedleft\arraybackslash}p{(\linewidth - 8\tabcolsep) * \real{0.2000}}
  >{\raggedleft\arraybackslash}p{(\linewidth - 8\tabcolsep) * \real{0.2000}}
  >{\raggedleft\arraybackslash}p{(\linewidth - 8\tabcolsep) * \real{0.2000}}
  >{\raggedleft\arraybackslash}p{(\linewidth - 8\tabcolsep) * \real{0.2000}}
  >{\raggedleft\arraybackslash}p{(\linewidth - 8\tabcolsep) * \real{0.2000}}@{}}
\toprule\noalign{}
\begin{minipage}[b]{\linewidth}\raggedleft
Particle streams
\end{minipage} & \begin{minipage}[b]{\linewidth}\raggedleft
Effective total tracks
\end{minipage} & \begin{minipage}[b]{\linewidth}\raggedleft
Trapped
\end{minipage} & \begin{minipage}[b]{\linewidth}\raggedleft
Incomplete
\end{minipage} & \begin{minipage}[b]{\linewidth}\raggedleft
Escaped
\end{minipage} \\
\midrule\noalign{}
\endhead
\bottomrule\noalign{}
\endlastfoot
100 & 500 & 159 & 341 & 0 \\
200 & 1000 & 316 & 684 & 0 \\
500 & 2500 & 824 & 1676 & 0 \\
1000 & 5000 & 1711 & 3289 & 0 \\
\end{longtable}
}

Percentage summary:

{\def\LTcaptype{none} % do not increment counter
\begin{longtable}[]{@{}rrr@{}}
\toprule\noalign{}
Total tracks & Trapped & Incomplete \\
\midrule\noalign{}
\endhead
\bottomrule\noalign{}
\endlastfoot
500 & 31.8\% & 68.2\% \\
1000 & 31.6\% & 68.4\% \\
2500 & 33.0\% & 67.0\% \\
5000 & 34.2\% & 65.8\% \\
\end{longtable}
}

Increasing total stochastic tracks from 500 to 5000 improved sampling
resolution, but it did not remove the incomplete-track problem. The
trapped fraction stayed around \textbf{31.6-34.2\%}, while incomplete
tracks remained dominant at around \textbf{65.8-68.4\%}.

For this reason, the DPM results are currently treated as
\textbf{diagnostic only}, not as final separator-efficiency evidence.

\begin{center}\rule{0.5\linewidth}{0.5pt}\end{center}

\subsection{3. What Was Found Out}\label{what-was-found-out}

\subsubsection{3.1 Main inlet-condition
finding}\label{main-inlet-condition-finding}

The inlet condition is close to the Purnanto 1600 kJ/kg case, but it is
not an exact match.

The reason is that the current simulation used:

\[
V = 26.81\ \text{m/s}
\]

rather than recalculating the inlet velocity from the current actual
inlet area.

The current inlet mass flow is:

\[
196.23\ \text{kg/s}
\]

whereas the Purnanto target is:

\[
197.61\ \text{kg/s}
\]

This gives a difference of approximately:

\[
-0.70\%
\]

This is small, but it should still be stated clearly so that the setup
is not presented as an exact mass-flow-matched Purnanto replication.

\begin{center}\rule{0.5\linewidth}{0.5pt}\end{center}

\subsubsection{3.2 Main outlet-flux
finding}\label{main-outlet-flux-finding}

The current outlet fluxes suggest:

{\def\LTcaptype{none} % do not increment counter
\begin{longtable}[]{@{}lr@{}}
\toprule\noalign{}
Indicator & Value \\
\midrule\noalign{}
\endhead
\bottomrule\noalign{}
\endlastfoot
Liquid separation indicator & 2.16\% \\
Steam outlet dryness & 97.02\% \\
\end{longtable}
}

The steam outlet appears relatively dry by mass. However, the liquid
outlet flow is very low, so the flux-based liquid separation value
should not be treated as a final separator efficiency yet.

This needs to be checked against:

\begin{itemize}
\tightlist
\item
  liquid volume fraction contours,
\item
  steam volume fraction contours,
\item
  velocity vectors,
\item
  pressure contours,
\item
  pathlines or streamlines,
\item
  residual convergence,
\item
  liquid accumulation inside the separator,
\item
  and outlet boundary behaviour.
\end{itemize}

\begin{center}\rule{0.5\linewidth}{0.5pt}\end{center}

\subsubsection{3.3 Visual evidence to add}\label{visual-evidence-to-add}

\paragraph{Liquid volume fraction
contours}\label{liquid-volume-fraction-contours}

\begin{quote}
Insert figure here.
\end{quote}

Comment on:

\begin{itemize}
\tightlist
\item
  where the liquid accumulates,
\item
  whether the liquid moves toward the wall,
\item
  whether liquid remains trapped in the separator,
\item
  and whether liquid approaches the steam outlet.
\end{itemize}

\paragraph{Steam volume fraction
contours}\label{steam-volume-fraction-contours}

\begin{quote}
Insert figure here.
\end{quote}

Comment on:

\begin{itemize}
\tightlist
\item
  whether a steam-rich core forms,
\item
  whether the steam outlet is mostly vapour,
\item
  and whether this supports the calculated outlet dryness.
\end{itemize}

\paragraph{Velocity contours / vectors}\label{velocity-contours-vectors}

\begin{quote}
Insert figure here.
\end{quote}

Comment on:

\begin{itemize}
\tightlist
\item
  whether the inlet creates strong swirl,
\item
  whether the flow circulates as expected,
\item
  whether stagnant or recirculating zones appear,
\item
  and whether there is any short-circuiting toward the steam outlet.
\end{itemize}

\paragraph{Pressure contours}\label{pressure-contours}

\begin{quote}
Insert figure here.
\end{quote}

Comment on:

\begin{itemize}
\tightlist
\item
  whether the pressure field is physically reasonable,
\item
  whether the inlet-to-outlet pressure drop is sensible,
\item
  and whether the pressure field supports cyclone-like flow behaviour.
\end{itemize}

\paragraph{Pathlines / streamlines}\label{pathlines-streamlines}

\begin{quote}
Insert figure here.
\end{quote}

Comment on:

\begin{itemize}
\tightlist
\item
  whether the flow spirals around the separator,
\item
  whether droplets or liquid-rich flow move downward,
\item
  whether some flow escapes upward,
\item
  and whether the pathlines explain the low liquid outlet flux.
\end{itemize}

\begin{center}\rule{0.5\linewidth}{0.5pt}\end{center}

\subsection{4. Meeting Summary}\label{meeting-summary}

The main point to explain in the meeting is that this simulation used
the \textbf{reported Purnanto velocity of 26.81 m/s}, not an
actual-area-corrected velocity. Because of this, the current inlet mass
flow is close to the Purnanto target but not exactly equal to it.

The current simulation gives a total inlet flow of approximately
\textbf{196.23 kg/s}, compared with the Purnanto target of
\textbf{197.61 kg/s}, which is about \textbf{0.70\% lower}.

The outlet fluxes suggest a relatively dry steam outlet, with
approximately \textbf{97.02\% outlet dryness}, but the liquid outlet
flow is very low, giving a flux-based liquid separation indicator of
only \textbf{2.16\%}. This should be treated as a preliminary result
until the contours, velocity field, pressure field, pathlines,
convergence, and outlet setup are checked.

For now, the DPM results should be treated as diagnostic only because
incomplete tracks remain high at around \textbf{66--68\%}.

\end{document}
